\section{Parcial 2025 Q2}

% =====================================================================
\subsection{Ejercicio 1 --- Modelo de seguridad para Palantir (CSO)}
% =====================================================================

\begin{consigna}
Palantir es una empresa startup que provee tecnología para control y mitigación de
\emph{swarms} de drones (entre otras cosas). Palantir necesita rediseñar todo su sistema
de seguridad para poder continuar trabajando con diversos ejércitos del mundo (a quienes
provee de tecnología \emph{top secret}). Ud, Ingeniera, Ingeniero, es el nuevo CSO.
\begin{enumerate}
  \item[a)] Detallar en este escenario, como CSO de Palantir, qué modelo de seguridad
  establecerían, qué políticas de control de acceso y qué reglas de seguridad.
  \item[b)] Describir cuál es la diferencia entre modelo de seguridad, política de control
  de acceso y reglas de seguridad.
\end{enumerate}
\end{consigna}

\paragraph{Temas.}
Clase 6 --- Políticas y Control de Acceso (\S Política de seguridad; \S ¿Qué es un modelo
de seguridad?; \S Bell--LaPadula; \S MAC vs.\ DAC; \S RBAC/ABAC/RuBAC).

\paragraph{Qué necesitás saber.}
\begin{itemize}
  \item \textbf{Modelo de seguridad:} modelo formal y lógico que define las reglas e
  interacciones que permiten implementar un mecanismo de control de acceso; con sustento
  matemático se demuestra que el sistema es seguro. Están especializados (confidencialidad,
  integridad, disponibilidad) y \emph{no alcanzan por sí solos} para todo el sistema: se
  aplican a partes.
  \item \textbf{Bell--LaPadula} (confidencialidad militar): clasificaciones $\{$TS, S, C, U$\}$
  con compartimentos formando un reticulado. Reglas: \emph{Read Down} (seguridad simple,
  $\Lvl(s)\dom\Lvl(o)$), \emph{Write Up} (propiedad $\ast$, $\Lvl(o)\dom\Lvl(s)$) y
  $\ast$-fuerte (leer+escribir $\Rightarrow$ misma clasificación). Es \textbf{control de
  flujo}: la información ``sube'', no ``baja''.
  \item \textbf{Política de control de acceso:} la definición de qué estados consideramos
  seguros (una ``foto en el tiempo'' que caracteriza si el estado actual es seguro). Se
  implementa con mecanismos \textbf{MAC} (mandatorio, basado en reglas universales del
  sistema, p.\ ej.\ el lattice de BLP) y/o \textbf{DAC} (discrecional, basado en identidad,
  el dueño/admin asigna permisos). Por comprensión: \textbf{RBAC} (roles), ABAC, RuBAC.
  \item \textbf{Reglas de seguridad:} los enunciados concretos que el mecanismo evalúa
  (las reglas de BLP, las entradas de una ACL, las reglas ordenadas tipo firewall).
\end{itemize}

\paragraph{Notas y conexiones.}
El escenario (datos \emph{top secret}, múltiples ejércitos que no deben verse entre sí)
es el caso canónico de \textbf{Bell--LaPadula}: confidencialidad + compartimentos
(\emph{need-to-know}). La idea de aislar proyectos con clientes distintos
es exactamente la motivación de los compartimentos del reticulado. BLP cubre la parte
mandatoria; el día a día (quién accede a qué bucket/repositorio) se gestiona con DAC/RBAC,
que tiene discrecionalidad fuerte (alguien asigna roles). Conecta con Clase 8 (separación
de privilegios, mínimo privilegio) y Clase 10 (Defense-in-depth, ZTA).

\paragraph{Resolución.}
\textbf{a)} Como CSO de Palantir propondría:
\begin{itemize}
  \item \textbf{Modelo de seguridad:} \textbf{Bell--LaPadula} como núcleo, por ser un
  modelo de \emph{confidencialidad} pensado para el ámbito militar y manejo de información
  clasificada. Se definen niveles de clasificación (Top Secret / Secret / Confidential /
  Unclassified) y, sobre todo, \textbf{compartimentos por cliente/proyecto} (cada ejército
  es una categoría incomparable con las demás), formando un reticulado. Donde BLP no alcance
  (sistemas de propósito general internos), se complementa con DAC/RBAC. Si además importa
  la \emph{integridad} de la información de control de drones, se puede sumar Biba o
  Clark--Wilson sobre esos componentes.
  \item \textbf{Políticas de control de acceso:} mandatorias (MAC) para la información
  clasificada --- las define el sistema y se aplican obligatoriamente a todos los accesos,
  cambios de atributos e intercambios entre sujetos; y discrecionales (DAC, típicamente
  \textbf{RBAC} por roles: analista, operador de drones, administrador, cliente militar)
  para la gestión cotidiana. \emph{Need-to-know} entre compartimentos.
  \item \textbf{Reglas de seguridad concretas:} las tres reglas de BLP ---
  \emph{Read Down} ($\Lvl(s)\dom\Lvl(o)$ para leer), \emph{Write Up} ($\Lvl(o)\dom\Lvl(s)$
  para escribir) y $\ast$-fuerte (lectura+escritura solo en la misma clasificación) ---
  más la dominancia con compartimentos: $(L,C)\dom(L',C')\iff L'\le L \wedge C'\subseteq C$.
  Dos clientes en compartimentos incomparables no pueden leerse ni escribirse mutuamente.
\end{itemize}

\textbf{b)} La diferencia (de lo más abstracto a lo más concreto):
\begin{itemize}
  \item El \textbf{modelo de seguridad} es el marco \emph{formal/matemático} (p.\ ej.\
  Bell--LaPadula): generaliza, demuestra propiedades y dice \emph{qué} hay que garantizar,
  sin atarse a una implementación.
  \item La \textbf{política de control de acceso} es la \emph{definición de qué estados son
  seguros} en este sistema concreto: a qué clasifico cada activo, qué sujetos existen, qué
  pueden hacer --- el ``qué se permite''. Decide entre mandatorio (MAC) y discrecional (DAC).
  \item Las \textbf{reglas de seguridad} son los enunciados concretos que el mecanismo
  evalúa para decidir cada acceso (las reglas de BLP, las entradas de una ACL, las reglas
  ordenadas de un firewall) --- el ``cómo se decide''.
\end{itemize}
En síntesis: el modelo \emph{fundamenta} la política, y la política se \emph{expresa}
mediante reglas que el mecanismo aplica.

% =====================================================================
\subsection{Ejercicio 2 --- Opción múltiple (pentesting, ley 25.326, buffer overflow)}
% =====================================================================

\begin{consigna}
Responder la opción más correcta.
\begin{itemize}
  \item ¿Cuáles son los pasos para implementar un pentesting exitoso?
  \begin{itemize}
    \item Recolección de información, Prueba de la seguridad, Generalización, Eliminación.
    \item Recolección de información, Prueba formal, Generalización, Eliminación.
    \item Recolección de información, Planteo de Hipótesis de falla, Prueba de la
    vulnerabilidad, Generalización, Eliminación.
    \item Planteo de Hipótesis de falla, Prueba de la vulnerabilidad, Verificación formal.
  \end{itemize}
  \item La ley argentina 25.326 de Protección de datos personales establece mecanismos y
  normativa para regular
  \begin{itemize}
    \item El Derecho al olvido:
    \item Las bases de datos públicas y privadas
    \item El manejo de las cookies en los navegadores de internet.
    \item Las historias clínicas.
  \end{itemize}
  \item ¿Cuál de las siguientes estrategias no tendría sentido en un ataque buffer overflow?
  \begin{itemize}
    \item Usar RUST
    \item Implementar un garbage collector
    \item Validar todos los datos de entrada.
    \item Implementar un Stack Canary.
  \end{itemize}
\end{itemize}
\end{consigna}

\paragraph{Temas.}
Clase 10b --- Pentesting (\S Metodología de hipótesis de falla); Clase 11 --- Protección
de Datos (\S Ley 25.326); Clase 8 --- Principios de Diseño y Vulnerabilidades
(\S CWE-787 / buffer overflow).

\paragraph{Qué necesitás saber.}
\begin{itemize}
  \item \textbf{Pasos del pentesting (de memoria, en orden):} Recolección de información
  $\rightarrow$ Planteo de hipótesis de falla $\rightarrow$ Prueba (de la vulnerabilidad)
  $\rightarrow$ Generalización $\rightarrow$ Eliminación (opcional). Equivocar el orden o
  los pasos penaliza. Notar que es \emph{hipótesis de falla}, no ``prueba formal'': la
  verificación formal es lo \emph{contrario} del pentesting.
  \item \textbf{Ley 25.326 (argentina):} regula las bases de datos \textbf{públicas y
  privadas} (hábeas data: ver/rectificar/borrar tus datos). El \textbf{Derecho al Olvido}
  lo regula el \textbf{GDPR europeo}, NO la ley argentina. La ley no especifica tecnología.
  \item \textbf{Buffer overflow (CWE-787, out-of-bounds write):} escribir fuera de los
  límites del buffer; en el stack puede pisar la dirección de retorno. Estrategias de
  defensa válidas: lenguajes que eliminan la aritmética de punteros / acceso a memoria
  seguro (RUST, Java/C\#), \emph{stack canary}, validar la entrada. Un \textbf{garbage
  collector} gestiona el \emph{ciclo de vida de memoria dinámica} (fugas, use-after-free),
  no impide escrituras fuera de límite en el stack.
\end{itemize}

\paragraph{Notas y conexiones.}
La primera y la última opción son las dos ``trampas'' que el profesor remarca: confundir
``Planteo de hipótesis de falla'' con ``prueba formal/de la seguridad'' (Clase 10b), y la
distinción ley argentina vs.\ GDPR (Clase 11), la más preguntable de protección de datos.
El buffer overflow conecta con Clase 8 (injection $\approx 90\%$; lenguajes ``seguros''
que mataron familias enteras de bugs) y con el principio de \emph{validar la entrada}.

\paragraph{Resolución.}
\begin{itemize}
  \item \textbf{Pasos del pentesting:} \textbf{``Recolección de información, Planteo de
  Hipótesis de falla, Prueba de la vulnerabilidad, Generalización, Eliminación''} (la 3.ª
  opción). Las otras sustituyen el paso clave por ``Prueba de la seguridad'' / ``Prueba
  formal'' o lo reemplazan por ``Verificación formal'' (que es la técnica opuesta).
  \item \textbf{Ley 25.326:} \textbf{``Las bases de datos públicas y privadas''}. El
  Derecho al olvido es GDPR; cookies (consentimiento) también es más bien GDPR; las
  historias clínicas son un caso particular (HIPAA en EE.UU.), no el objeto de la 25.326.
  \item \textbf{Buffer overflow:} la estrategia que \textbf{no tiene sentido} es
  \textbf{``Implementar un garbage collector''}: un GC administra memoria dinámica, no
  evita la escritura fuera de los límites de un buffer (el desbordamiento clásico). Usar
  RUST (acceso a memoria seguro), validar la entrada y un stack canary sí mitigan.
\end{itemize}

% =====================================================================
\subsection{Ejercicio 3 --- Encripción homomórfica}
% =====================================================================

\begin{consigna}
La encripción homomórfica aparece como un escenario donde es posible almacenar datos en la
nube directamente encriptados y donde se pueden realizar operaciones de cálculo directamente
sobre los datos, obteniendo un resultado que al desencriptarlo coincidiría con el resultado
obtenido de la operación.

Por ejemplo: $f(Enc_k(x)) = Enc_k(f(x))$
\begin{enumerate}
  \item[a)] Explicar qué escenario de protección de datos podría ser beneficioso este esquema.
\end{enumerate}
\end{consigna}

\paragraph{Temas.}
Clase 11 --- Protección de Datos (\S Encripción homomórfica; \S Estados del dato:
\emph{in-use}).

\paragraph{Qué necesitás saber.}
La encripción homomórfica permite \textbf{operar sobre los datos cifrados sin descifrarlos}:
$f(\mathrm{Enc}_k(x)) = \mathrm{Enc}_k(f(x))$. Es decir, operar sobre el texto cifrado y
luego descifrar da el mismo resultado que operar sobre el texto plano. Ataca el estado más
difícil de proteger: el dato \emph{en uso} (in-use), cuando normalmente hay que descifrarlo
para procesarlo.

\paragraph{Notas y conexiones.}
El control de acceso no alcanza para proteger datos: el dato fluye (RR.HH.\ $\to$ App $\to$
nube). Cuando delegamos cómputo a un tercero (la nube), el problema es que para procesar hay
que exponer el dato en claro. La homomórfica rompe ese compromiso. Conecta con los estados
del dato (Clase 11) y con \emph{shared tenancy} / confidencialidad en la nube (Clase 9).

\paragraph{Resolución.}
El escenario beneficiado es la \textbf{computación delegada en la nube preservando la
privacidad}: un tercero (proveedor cloud) procesa, analiza o agrega datos sensibles
\textbf{sin verlos nunca en claro}. El cliente cifra ($\mathrm{Enc}_k(x)$), envía el cifrado,
el servidor computa $f$ sobre el cifrado y devuelve $\mathrm{Enc}_k(f(x))$; solo el cliente,
con su clave, descifra el resultado:
\begin{equation*}
\mathrm{Dec}_k\big(\mathrm{Enc}_k(f(x))\big) = f(x).
\end{equation*}
Ejemplos concretos: análisis estadístico/ML sobre datos médicos o financieros tercerizado a
un proveedor que no debe acceder a los datos en claro (HIPAA, datos de salud anonimizados),
votaciones electrónicas, o consultas a una base de datos en la nube sin revelar ni la consulta
ni los registros. Resuelve la protección del dato \emph{en uso}, que de otro modo obligaría a
descifrar para procesar.

% =====================================================================
\subsection{Ejercicio 4 --- Secreto compartido de Shamir}
% =====================================================================

\begin{consigna}
Considerar un esquema de secreto compartido de Shamir $(k,n)=(3,5)$ en $\mathbb{Z}_7$, con el
conjunto de sombras
\begin{equation*}
C = \{(1,2),\ (2,2),\ (3,3),\ (4,5),\ (5,1)\}.
\end{equation*}
\begin{enumerate}
  \item[a)] Recuperar el secreto.
  \item[b)] Cuál es el polinomio generador.
  \item[c)] Considerar el anillo de polinomios $\mathbb{F}_2[x]$ y el campo de extensión
  \begin{equation*}
  \mathbb{F}_{2^3} \cong \mathbb{F}_2[x]/(p(x)),\qquad p(x)=x^3+x+1,
  \end{equation*}
  denotando $\alpha$ la clase de $x$ en $\mathbb{F}_{2^3}$ (es decir $\alpha = x \bmod p(x)$).
  Se define un esquema de Shamir $(k,n)=(3,4)$ sobre $\mathbb{F}_{2^3}$. Las sombras
  entregadas son
  \begin{equation*}
  C=\{(1,\ \alpha+1),\ (\alpha,\ 1+\alpha+\alpha^2),\ (\alpha+1,\ 1),\ (\alpha^2,\ \alpha+1)\},
  \end{equation*}
  donde los puntos se escriben como $(t,s_t)$ con $t\in\mathbb{F}_{2^3}$ y
  $s_t\in\mathbb{F}_{2^3}$. Detallar por qué razón el algoritmo secreto compartido de Shamir
  puede también plantearse de esta manera.
\end{enumerate}
\end{consigna}

\paragraph{Temas.}
Clase 8 --- Principios de Diseño y Vulnerabilidades (\S El secreto compartido de Shamir, en
términos de entropía); flujo de información (entropía). Mecánica: interpolación de Lagrange
sobre un cuerpo finito.

\paragraph{Qué necesitás saber.}
\begin{itemize}
  \item En un esquema $(T,N)$ de Shamir hay $N$ sombras y se necesitan \textbf{al menos $T$}
  para recuperar el secreto $S$. El secreto se codifica como el término independiente de un
  polinomio de grado $T-1$; cada sombra es un punto $(x_i, p(x_i))$. Con $<T$ puntos la
  entropía del secreto \emph{no cambia} (no hay flujo); con $\ge T$, $H(S)\to 0$.
  \item Aquí $T=3$, así que $p$ tiene grado $2$: $p(x)=a_0+a_1x+a_2x^2$, y el secreto es
  $S=a_0=p(0)$. Se recupera con \textbf{interpolación de Lagrange} usando 3 puntos, toda la
  aritmética módulo $7$ (cuerpo $\mathbb{Z}_7$).
  \item La validez del esquema (parte c) depende solo de que el conjunto de coeficientes /
  evaluaciones viva en un \textbf{cuerpo}: lo que se necesita es poder sumar, restar,
  multiplicar y \textbf{dividir} (existencia de inversos) para que Lagrange funcione.
\end{itemize}

\paragraph{Notas y conexiones.}
Shamir aparece en Clase 8 como ejemplo de flujo de información medido con entropía: ``con
menos de $T$ sombras no hay flujo'' equivale a ``$H(S)$ no cambia''. Aquí se pide además la
mecánica (Lagrange). La parte c) conecta con la Primera Parte (cuerpos finitos
$\mathbb{F}_{2^n}$, aritmética polinomial módulo un irreducible, igual que en AES).

\paragraph{Resolución.}

\textbf{a) y b)} Buscamos $p(x)=a_0+a_1x+a_2x^2$ en $\mathbb{Z}_7$. Con tres sombras, por
ejemplo $(1,2),(2,2),(3,3)$, planteamos el sistema:
\begin{align*}
a_0 + a_1 + a_2 &\equiv 2 \pmod 7,\\
a_0 + 2a_1 + 4a_2 &\equiv 2 \pmod 7,\\
a_0 + 3a_1 + 2a_2 &\equiv 3 \pmod 7.
\end{align*}
Resolviendo (eliminación de Gauss módulo $7$, recordando que en $\mathbb{Z}_7$ los inversos
son $2^{-1}=4$, $3^{-1}=5$, $4^{-1}=2$, $5^{-1}=3$, $6^{-1}=6$) se obtiene
\begin{equation*}
a_0 = 3,\qquad a_1 = 2,\qquad a_2 = 4.
\end{equation*}
Luego el \textbf{polinomio generador} es
\begin{equation*}
\boxed{\,p(x) = 4x^2 + 2x + 3 \pmod 7\,}
\end{equation*}
y el \textbf{secreto} es el término independiente:
\begin{equation*}
\boxed{\,S = p(0) = a_0 = 3\,.}
\end{equation*}
\emph{Verificación} (consistencia del esquema): $p(1)=4+2+3=9\equiv 2$, $p(2)=16+4+3=23\equiv
2$, $p(3)=36+6+3=45\equiv 3$, $p(4)=64+8+3=75\equiv 5$, $p(5)=100+10+3=113\equiv 1$
$\pmod 7$. Las cinco sombras caen sobre el polinomio, así que el reparto es válido y
cualquier subconjunto de $3$ sombras recupera el mismo secreto.

\textbf{c)} El algoritmo de Shamir \textbf{no depende de que trabajemos en $\mathbb{Z}_p$
con $p$ primo}: lo único que usa es que las coordenadas vivan en un \textbf{cuerpo (campo)}.
La construcción y la recuperación se basan en la \textbf{interpolación de Lagrange}, que
requiere sumar, restar, multiplicar y \textbf{dividir} --- y dividir solo es posible si todo
elemento no nulo tiene inverso multiplicativo, que es justamente la propiedad que define a un
cuerpo.

$\mathbb{F}_{2^3}\cong\mathbb{F}_2[x]/(p(x))$ con $p(x)=x^3+x+1$ \textbf{irreducible} sobre
$\mathbb{F}_2$ es un cuerpo finito de $2^3=8$ elementos (las clases
$\{0,1,\alpha,\alpha+1,\alpha^2,\dots\}$), con suma y producto de polinomios módulo $p(x)$ y,
por ser cuerpo, \textbf{inversos para todo elemento no nulo}. Por lo tanto:
\begin{itemize}
  \item Se puede definir un polinomio secreto de grado $T-1=2$ con coeficientes en
  $\mathbb{F}_{2^3}$, evaluarlo en $N=4$ puntos distintos $t\in\mathbb{F}_{2^3}$ para generar
  las sombras $(t,s_t)$, y
  \item recuperar el secreto $S=p(0)$ con la misma fórmula de Lagrange, ahora con la
  aritmética de $\mathbb{F}_{2^3}$ (los inversos se calculan módulo $p(x)$).
\end{itemize}
Es decir, Shamir se plantea idénticamente sobre cualquier cuerpo finito $\mathbb{F}_q$ (sea
$q$ primo como $\mathbb{Z}_7$, o potencia de primo como $\mathbb{F}_{2^3}$); usar
$\mathbb{F}_{2^n}$ es natural cuando los secretos son cadenas de bits, exactamente como en
AES. \textbf{TODO:} verificar numéricamente la recuperación en $\mathbb{F}_{2^3}$ no se pide
explícitamente; la consigna solo pide \emph{justificar por qué} es válido plantearlo así
(la respuesta es: porque $\mathbb{F}_{2^3}$ es un cuerpo y Lagrange solo necesita un cuerpo).

% =====================================================================
\subsection{Ejercicio 5 --- ``Garantizar la seguridad'' con pentesting (Tesla)}
% =====================================================================

\begin{consigna}
Tesla implementó un sistema de mejoramiento de la calidad de software nuevo. Su director,
Charles Kharpaty, anunció que estarían utilizando un nuevo esquema que garantizaba la
seguridad basado en estrictas pruebas de pentesting.
\begin{enumerate}
  \item[a)] Provea argumentos sólidos para deslizarle al oído de Elon para que lo despida o
  lo promueva a CSO.
\end{enumerate}
\end{consigna}

\paragraph{Temas.}
Clase 10b --- Pentesting (\S Limitaciones: el pentesting NO garantiza la seguridad;
\S Verificación formal vs.\ pentesting).

\paragraph{Qué necesitás saber.}
\textbf{Gancho/trampa de parcial:} el pentesting es una buena estrategia pero \textbf{NO
garantiza la seguridad}. ``Garantizar la seguridad es algo que no se puede hacer jamás''
(palabras del profesor). Si un enunciado afirma que algo ``garantiza la seguridad'' es
incorrecto, y \textbf{hay penalización si se les escapa}. Razones:
\begin{itemize}
  \item El pentesting prueba la \textbf{existencia} de vulnerabilidades, \textbf{nunca su
  ausencia} (Dijkstra: el testing muestra presencia de bugs, no su ausencia). Quien sí
  prueba ausencia es la \emph{verificación formal}, pero solo en un algoritmo/ambiente
  acotado incluyendo todos los factores --- y \emph{tampoco} da garantía total para un
  sistema completo.
  \item No es sistemático: depende de la capacidad del tester para formular hipótesis.
  \item Cada test vale para \emph{ese} sistema en \emph{ese} momento (si cambia la API/v3,
  las pruebas quedan obsoletas).
  \item No reemplaza un buen diseño, especificación, implementación y testing; se suma
  \emph{después}.
\end{itemize}

\paragraph{Notas y conexiones.}
Es el mismo gancho que el Ejercicio 2 (``Prueba formal'' vs.\ hipótesis de falla) visto desde
el lado del enunciado-trampa. Conecta con Clase 6/8 (la seguridad es un estado que se debe
mantener; un buen diseño y los principios de Saltzer--Schroeder son la base) y con Clase 10
(DevSecOps: el pentesting es una capa más, no la garantía).

\paragraph{Resolución.}
\textbf{Despedirlo} (o, como mínimo, corregir el anuncio). El error de fondo es la afirmación
de que el esquema \textbf{``garantiza la seguridad''}: \emph{garantizar la seguridad es
imposible}. Argumentos sólidos para el oído de Elon:
\begin{enumerate}
  \item \textbf{El pentesting nunca prueba ausencia de vulnerabilidades, solo existencia.}
  Encontrar fallas demuestra que las hay; \emph{no} encontrarlas no demuestra que no las
  haya (Dijkstra). Anunciar ``garantía de seguridad'' es técnicamente falso.
  \item \textbf{No es sistemático:} su calidad depende de la habilidad del tester para
  plantear hipótesis de falla; dos testers distintos cubren cosas distintas.
  \item \textbf{Es válido para ese sistema en ese momento:} cualquier cambio (nueva versión,
  nueva integración) invalida buena parte de las pruebas. No hay garantía persistente.
  \item \textbf{No reemplaza el buen diseño:} la seguridad se construye con principios de
  diseño, especificación, implementación cuidada y cobertura de testing; el pentesting se
  agrega \emph{después} para comprobar, no para garantizar.
  \item Si se quisiera la mayor garantía posible, lo más cercano sería la \textbf{verificación
  formal} (prueba ausencia, pero solo en un dominio acotado y aun así no total), no el
  pentesting.
\end{enumerate}
Conclusión: Kharpaty merece ser \textbf{despedido (o corregido) por prometer una garantía que
no existe}; un buen CSO diría ``reducimos el riesgo y aumentamos la confianza'', nunca
``garantizamos la seguridad''.
